% ESP Paper Template README
% ============================

% SETUP CHECKLIST  
% ------------------
% 0. Make a copy of this template File > Make a copy 
% 1. Import: 
%     - Bibliography (either upload your own, or replace the content in references.bib) 
%     - Figures 
%     - Logos of external collaborators (if applicable) 
% 2. Fill in your title, authors, and affiliations at the top of main.tex
% 3. Replace placeholder sections with your content
% 4. Compile!

% In case you have compiling errors, double-check that you are compiling using XeLaTeX. Check under File > Settings > Compiler > XeLaTeX

 
% ESP-ONLY VS. EXTERNAL COLLABORATIONS
% -------------------------------------- 

%   ESP-only:  
%     - Use \documentclass{esp-paper} at the top of main.tex 
%     - Single \affil{Earth Species Project} for all authors
%     - Shows only the ESP logo on Page 1 and subsequent headers
      
%   External:  
%     - Use \documentclass[external]{esp-paper} at the top of main.tex 
%     - Numbered superscripts ($^1$, $^2$) per author and affiliation
%     - Including collaborator's logos (as it makes sense, use your best judgment if there is a very long list of external affiliations) 
 
% CITATIONS
% -------------------------------------- 
%   - Use \citep{key} for parenthetical citations: (Smith, 2023)
%   - Use \citet{key} for inline citations: Smith (2023)
%   - When copying over from another .tex file that used \cite{}, the parenthetical formatting may not transfer over. Do a Find + Replace to ensure you are using \citep{}.  


% OPTIONAL FEATURES
% -----------------
%   Short title:  
%     - Shown in the running header from page 2 onward
%     - If your main title is short enough, you could reuse it here
%     - You may omit if the main title is too long, and does not have a convenient shorter title  
  
%   Two-column layout:
%     - If you need 2 columns, uncomment out the lines 103 and 185
%     - \begin{multicols}{2} ... \end{multicols} 
  
%   Collaborator logos: 
%     - Add images to logos/ and uncomment the \extralogos{} block
    
%   Custom commands:    
%     - Add shorthands in the CUSTOM COMMANDS section of main.tex 